\chapter{Clustering}\label{clustering}

\textbf{Learning Objectives}

\emph{Overview starts}

\begin{itemize}
\tightlist
\item
  Identifying and understanding the main clustering algorithms.
\item
  Exploring practical examples of different clustering methods.
\item
  Understanding the techniques for clustering text or documents to uncover
  latent patterns.
\item
  Analyzing and extracting information from clustering to aid decision-making.
\item
  Experimenting with clustering and apply it in different
  contexts.
\end{itemize}

\emph{Overview ends}

\section{Grouping text and linguistic
elements}\label{grouping-text-and-linguistic-elements}

Humans tend to categorize (Harnad 2017): This is a cognitive
shortcut which identifies commonalities and creates groups, enabling the efficient processing of
information. This is part of human cognition and
allows us to act in a predetermined way when a specific, already familiar, situation
arises. It is also serves as a method of evaluating homogeneity in data so that we can make decisions based on a restricted number of
groups with specific characteristics. It is similar to summarizing: We
condense a lot of information into few important key points that
describe the whole in a concise manner.

One practical application of categorizing is clustering. This
involves identifying shared characteristics within data and grouping
similar items accordingly. There are several clustering
techniques, as well as many fields of application. In marketing,
clustering can also be referred to as segmentation, i.e., when we try to find a
segment of a population that exhibits similar behaviors and patterns. Of
clustering methods one is model-based
and uses statistical probability models, another is density-based
clustering, which uses a model that groups objects based on objective functions, while yet another is distance-based
and has a strong exploratory component. Various measures can be
used for this last version, including Hamming for categorical data,
Manhattan and Euclidean for continuous data, Jaccard for binary data,
Mahalanobis for correlated data, or cosine similarity for text data.
There are also hybrid methods which combine various
clustering methods (Vries, Olech, and Moscato 2019). 

A similar concept
to clustering is classification. This uses existing
classes with known features: New unclassified elements are
tested against these features and then assigned to the appropriate
class, so classification is a supervised method (Sen, Hajra, and Ghosh
2020). Clustering seeks to uncover these commonalities autonomously,
often through iterative, unsupervised mathematical operations within
algorithms. Often we do not know how many clusters we can have and
need to conduct several trials before we can extract meaningful insights into the data. This is a brief outline of two important concepts, clustering
and classification, and in this chapter we will concentrate on the principles and
application of the former.

In clustering, the goal is to extract latent groups, or clusters, that
contain a certain number of elements. These clusters may range from those which contain a single element to those in which all elements within a dataset are present. In any
clustering algorithm the concept of similarity is twofold: high
similarity (homogeneity) within clusters and marked dissimilarity
(heterogeneity) between different clusters, as illustrated in
Figure~\ref{fig-clusteringin}. When applying a clustering algorithm,
several factors are of interest, including the number of clusters
generated, their relative distances, sizes, shapes, and densities.

Text clustering gives us a quick overview of the content of a corpus by
extracting its most representative features. Grouping together
the words that are most closely related to each other (Struhl 2015) helps us to identify underlying topics and themes. These
methods can be very useful, for example, when clustering search results or
finding similar word groups in emails from customers. Such insights are
essential in refining segmentation strategies and discovering recurring
issues in customer communications, and so helping to develop one-to-one
marketing.

\begin{figure}

\centering{

\includegraphics[width=0.5\textwidth,height=\textheight]{figure/chapter5/cluster.png}

}

\caption{\label{fig-clusteringin}Cluster similarities \newline Authors'
own figure}

\end{figure}%

\section{Proximity and similarity}\label{proximity-and-similarity}

\subsection{Euclidean distance}\label{sec-euclidean-distance}

In text mining, clustering can be applied at various levels, for
entire books, chapters, paragraphs, sentences, or even individual words.
Clustering words means grouping terms that appear at similar
frequencies across different documents (see the definition of document in
the introductory chapter), while clustering documents allows us to
understand the similarities between them. Authorship attribution is one
practical application, where texts are grouped to assign them to
specific authors. Each author's unique writing style, manifested through
their use of stop words, distinctive expressions, word frequencies, and
punctuation, becomes an identifier for such analyses (Stamatatos et al.
2016). Research has also shown gender-based differences in writing
styles, indicating potential gender segmentation in texts (Cheng,
Chandramouli, and Subbalakshmi 2011).

Frequency (i.e.,~how many times a term appears in a text) allows us
to use algorithms that involve the concept of Euclidean distance (Huang
et al. 2008). Documents with a high degree of similarity of word
frequencies are considered closer in terms of Euclidean distance,
meaning that the distance between them is minimal. A simple
example will help to grasp the basic concept.

Consider an unlikely scenario of three documents, each containing just
two words. Each word's occurrence in a document is quantified by its
frequency, expressed as the number of times it appears. To
demonstrate this concept, we will use R to replicate the data presented
in Table~\ref{tbl-firsttable} where the numerical representation of our
example is shown.

\begin{longtable}[]{@{}lll@{}}
\caption{Three documents, two
words}\label{tbl-firsttable}\tabularnewline
\toprule\noalign{}
& bank & engine \\
\midrule\noalign{}
\endfirsthead
\toprule\noalign{}
& bank & engine \\
\midrule\noalign{}
\endhead
\bottomrule\noalign{}
\endlastfoot
doc1 & 20 & 11 \\
doc2 & 22 & 12 \\
doc3 & 12 & 27 \\
\end{longtable}

To begin our analysis, we define the variables \texttt{doc1},
\texttt{doc2} and \texttt{doc3} containing the frequency counts of two
specific words within the respective documents. The function
\texttt{rbind()} is used to merge the individual \texttt{doc} vectors
into a single data frame consolidating them in terms of row. It is worth noting
the existence of \texttt{cbind()}, a complementary function that
consolidates columns.

After binding these vectors, we use the function \texttt{rename()}
to assign meaningful column names to the newly formed table stored in the variable \texttt{docs}: \texttt{bank} and \texttt{engine}. To represent the co-occurrence frequencies of these two words visually across the documents, we construct a
scatterplot. The function \texttt{ggplot()} and its associated
attributes are used to display our data. The function
\texttt{geom\_point()} is applied to generate a scatterplot,
\texttt{labs()} and \texttt{geom\_label()} enhance the plot by adding labels
and creating an aesthetically pleasing layout, see
Figure~\ref{fig-euclideandist}.

\emph{Program code starts}

\begin{Shaded}
\begin{Highlighting}[]
\FunctionTok{library}\NormalTok{(tidyverse)}
\CommentTok{\# Example, the word bank (appears 20 times in the first document), }
\CommentTok{\# engine (appears 11 times in the same document)}
\NormalTok{doc1 }\OtherTok{\textless{}{-}} \FunctionTok{c}\NormalTok{(}\DecValTok{20}\NormalTok{, }\DecValTok{11}\NormalTok{)}
\NormalTok{doc2 }\OtherTok{\textless{}{-}} \FunctionTok{c}\NormalTok{(}\DecValTok{22}\NormalTok{, }\DecValTok{12}\NormalTok{)}
\NormalTok{doc3 }\OtherTok{\textless{}{-}} \FunctionTok{c}\NormalTok{(}\DecValTok{12}\NormalTok{, }\DecValTok{27}\NormalTok{)}

\NormalTok{docs }\OtherTok{\textless{}{-}} \FunctionTok{rbind}\NormalTok{(doc1, doc2, doc3) }\SpecialCharTok{\%\textgreater{}\%}
\NormalTok{  as.data.frame }\SpecialCharTok{\%\textgreater{}\%}
  \FunctionTok{rename}\NormalTok{(}\AttributeTok{bank=}\DecValTok{1}\NormalTok{, }\AttributeTok{engine=}\DecValTok{2}\NormalTok{)}
  
\FunctionTok{ggplot}\NormalTok{(docs, }\FunctionTok{aes}\NormalTok{(}\AttributeTok{x=}\NormalTok{bank, }\AttributeTok{y=}\NormalTok{engine)) }\SpecialCharTok{+} 
  \FunctionTok{geom\_point}\NormalTok{() }\SpecialCharTok{+} 
  \FunctionTok{labs}\NormalTok{(}\AttributeTok{x=}\StringTok{"bank"}\NormalTok{, }\AttributeTok{y=}\StringTok{"engine"}\NormalTok{) }\SpecialCharTok{+} 
  \FunctionTok{geom\_label}\NormalTok{(}\AttributeTok{label=}\FunctionTok{rownames}\NormalTok{(docs), }\AttributeTok{nudge\_x =} \FloatTok{0.35}\NormalTok{, }\AttributeTok{nudge\_y =} \FloatTok{0.35}\NormalTok{)}
\end{Highlighting}
\end{Shaded}

\begin{figure}[H]

\centering{

\includegraphics{5.Chapter_files/figure-pdf/fig-euclideandist-1.pdf}

}

\caption{\label{fig-euclideandist}Documents similitude \newline Authors'
own figure}

\end{figure}%

\emph{Program code ends}

After we bind the documents and the words together in a tabular
structure, the data frame, we can calculate the distance between
the documents. Note that the concept of distance is
associated with the concept of similarity: The smaller the distance, the
higher the similarity. Figure~\ref{fig-euclideandist} gives a
visual interpretation of this concept, illustrating that \texttt{doc1}
and \texttt{doc2} are closer to one another when compared with
\texttt{doc3}. Mathematically, this can be calculated using
Euclidean distance (Jockers and Thalken 2020b). On a
two-dimensional plane, this is analogous to the hypotenuse in the theorem of Pythagoras. Be aware that Eucledian distance is not the only
method that can be used to calculate the distance between two or more elements. The work of
Deza and Deza (2009) offers an extensive overview of distance metrics
covering a wide range of calculation methods.

The equation to calculate the distance between \texttt{doc1} and
\texttt{doc2} uses \(w\) to represent word frequencies in the
documents. It is essential to understand that \(1^{st}w_{doc_1}\),
\(1^{st}w_{doc_2}\), etc., represent the frequencies of the first words
in \textbf{doc1} and \textbf{doc2} respectively. The pattern continues
for all words, up the \(n^{th}\) word. This is a simplification which
assumes that both documents contain exactly the same set of words. For
words that are present in one document and absent in the other, their
frequency in the document where they are absent is considered to be zero. The equation
calculates the Euclidean distance by summing the squares of differences
in word frequencies for each word across the two documents and then
taking the square root of this sum.
\[dist(doc_1, doc_2)=\sqrt{(1^{st}w_{doc_1}-1^{st}w_{doc_2})^2+(2^{nd}w_{doc_1}-2^{nd}w_{doc_2})^2+\ldots+(n^{th}w_{doc_1}-n^{th}w_{doc_2})^2}\]

In our case, we do not have \emph{n} words, but only two: \emph{bank}
and \emph{engine}. Let us calculate the distance between the documents
and round the result to the first two decimals.

\[
dist(doc_1, doc_2)=\sqrt{(20-22)^2+(11-12)^2}=2.24 
\] \[
dist(doc_1, doc_3)=\sqrt{(20-12)^2+(11-27)^2}=17.89
\] \[
dist(doc_2, doc_3)=\sqrt{(22-12)^2+(12-27)^2}=18.03
\]

As was expected from the scatter plot, the Eucledian distance between
documents \(d_1\) and \(d_2\) is considerably smaller than the distance
between \(d_1\) and \(d_3\), or between \(d_2\) and \(d_3\). To
facilitate the calculation of these distances, especially in cases
involving many documents with multiple features, we can use the R
function \texttt{dist()}. This efficiently calculates the
distances of pairs between the documents.

\emph{Program code starts}

\begin{Shaded}
\begin{Highlighting}[]
\CommentTok{\# Compute pariwise distances between documents}
\FunctionTok{dist}\NormalTok{(docs)}
\end{Highlighting}
\end{Shaded}

\begin{verbatim}
          doc1      doc2
doc2  2.236068          
doc3 17.888544 18.027756
\end{verbatim}

\emph{Program code ends}

Euclidean distance is a simple method for calculating proximity, but
in marketing it is not used frequently due to its sensitivity to the
magnitude of the term frequencies. Another metric, called cosine
distance is preferred, as it measures the angle between vectors: We will
discuss this in more in detail in
Chapter~\ref{sec-extracting-information-from-embeddings}. In marketing
applications, cosine distance helps in the analysis of text similarity, such as
clustering customer reviews or segmenting product descriptions (Berger
et al. 2020).

\subsection{Cosine similarity}\label{cosine-similarity}

We can calculate the cosine similarity of two vectors by using the
data from the previous session. We can consider a document
as a vector, with the elements of the vector being the frequencies of the words
in a document (Han, Pei, and Tong 2022). The cosine similarity
between two vectors \(\vec{A}\) and \(\vec{B}\) is defined as:

\begin{equation}\phantomsection\label{eq-cosine_similarity}{
\text{sim}(\vec{A}, \vec{B}) = \frac{\sum_{i=1}^n A_i \cdot B_i}{\sqrt{\sum_{i=1}^n A_i^2} \cdot \sqrt{\sum_{i=1}^n B_i^2}}
}\end{equation}

In Equation~\ref{eq-cosine_similarity}, \(A_i\) and \(B_i\) are the
\(i\text{-th}\) components, in our case word frequencies, of the two
vectors. The numerator represents their dot product while the
denominator is the product of their Euclidean magnitudes. We have the
three documents represented as vectors, \(\vec{doc_1} = (20, 11)\),
\(\vec{doc_2} = (22, 12)\), \(\vec{doc_2} = (12, 27)\). The cosine
similarity between \(\vec{doc_1}\) and \(\vec{doc_2}\) is calculated
as

\[
\text{sim}(\vec{doc_1}, \vec{doc_2}) = \frac{(20 \cdot 22) + (11 \cdot 12)}{\sqrt{20^2 + 11^2} \cdot \sqrt{22^2 + 12^2}}
= \frac{572}{22.83 \cdot 25.07} \approx 0.998
\] Similarly, we can calculate the cosine similarity between
\(\vec{doc_1}\) and \(\vec{doc_3}\), and between \(\vec{doc_2}\) and
\(\vec{doc_2}\). The closer the similarity measure between the two vectors
is to \(1\), the more similar they are. In R we can use the package
\texttt{lsa} to speed up the computation.

\emph{Program code starts}

\begin{Shaded}
\begin{Highlighting}[]
\FunctionTok{library}\NormalTok{(lsa)}

\CommentTok{\# Example data: term frequencies}
\NormalTok{doc1 }\OtherTok{\textless{}{-}} \FunctionTok{c}\NormalTok{(}\DecValTok{20}\NormalTok{, }\DecValTok{11}\NormalTok{)}
\NormalTok{doc2 }\OtherTok{\textless{}{-}} \FunctionTok{c}\NormalTok{(}\DecValTok{22}\NormalTok{, }\DecValTok{12}\NormalTok{)}
\NormalTok{doc3 }\OtherTok{\textless{}{-}} \FunctionTok{c}\NormalTok{(}\DecValTok{12}\NormalTok{, }\DecValTok{27}\NormalTok{)}

\CommentTok{\# Combine into a matrix}
\NormalTok{docs }\OtherTok{\textless{}{-}} \FunctionTok{rbind}\NormalTok{(doc1, doc2, doc3)}

\CommentTok{\# Calculate cosine similarity matrix}
\NormalTok{cosine\_sim\_matrix }\OtherTok{\textless{}{-}} \FunctionTok{cosine}\NormalTok{(}\FunctionTok{t}\NormalTok{(docs))}

\CommentTok{\# View the cosine similarity matrix}
\NormalTok{cosine\_sim\_matrix}
\end{Highlighting}
\end{Shaded}

\begin{verbatim}
          doc1      doc2      doc3
doc1 1.0000000 0.9999939 0.7962479
doc2 0.9999939 1.0000000 0.7941278
doc3 0.7962479 0.7941278 1.0000000
\end{verbatim}

\emph{Program code ends}

It is easy to observe that \texttt{doc1} and \texttt{doc2} are the
most similar, while \texttt{doc3} is less.

\section{Hierarchical agglomerative
clustering}\label{hyerarchical-agglomerative-clustering}

If we do not want to cluster or find document similarities, but we need
to focus on the words, a directly related application, exemplifying
the principles discussed in Section~\ref{sec-euclidean-distance}, is
hierarchical clustering (HC), one of the simplest forms of clustering.
HC has two primary methodologies: agglomerative and divisive. In
agglomerative (or bottom-up) clustering, each element initially
represents a distinct cluster. Then, they are merged, or agglomerated,
into increasingly larger clusters, culminating in a single, comprehensive
cluster. In the divisive (or top-down) clustering, the algorithm begins
with one large cluster which is systematically divided into smaller and
smaller ones, ultimately reaching the level of individual elements with a
tree-like taxonomy. The analogy with a tree (dendrogram) is
straightforward: the end branches are individual objects, the nodes where they divide signify
the clusters, and the root contains the entire dataset, as illustrated in
Figure~\ref{fig-clusteringhac}.

\begin{figure}

\centering{

\includegraphics[width=0.99\textwidth,height=\textheight]{figure/chapter5/hac.png}

}

\caption{\label{fig-clusteringhac}Hierarchical Agglomerative/Divisive
clustering \newline Authors' own figure}

\end{figure}%

We have learned how to calculate the distance between objects using
Eucledian distance, but a cluster contains one or more than one element.
So, if the distance between individual elements is well understood, how
can we calculate the distance between clusters? There are several ways
to do this. The distance between two clusters can be given by the
distance between either the closest or furthest elements contained in the two clusters. Or it is possible to calculate the mean distance between
all the pairs of elements, one from each cluster. Both of these methods,
known collectively as linkage criteria, influence the shape and size of
the clusters formed. To better visualize these concepts,
Table~\ref{tbl-linkage} provides diagrams and succinct descriptions.

\begin{longtable}[]{@{}
  >{\raggedright\arraybackslash}p{(\columnwidth - 4\tabcolsep) * \real{0.1806}}
  >{\raggedright\arraybackslash}p{(\columnwidth - 4\tabcolsep) * \real{0.6389}}
  >{\raggedright\arraybackslash}p{(\columnwidth - 4\tabcolsep) * \real{0.1806}}@{}}
\caption{Linking Criteria}\label{tbl-linkage}\tabularnewline
\toprule\noalign{}
\begin{minipage}[b]{\linewidth}\raggedright
Linkage Type
\end{minipage} & \begin{minipage}[b]{\linewidth}\raggedright
Explanation
\end{minipage} & \begin{minipage}[b]{\linewidth}\raggedright
Diagram
\end{minipage} \\
\midrule\noalign{}
\endfirsthead
\toprule\noalign{}
\begin{minipage}[b]{\linewidth}\raggedright
Linkage Type
\end{minipage} & \begin{minipage}[b]{\linewidth}\raggedright
Explanation
\end{minipage} & \begin{minipage}[b]{\linewidth}\raggedright
Diagram
\end{minipage} \\
\midrule\noalign{}
\endhead
\bottomrule\noalign{}
\endlastfoot
Single Linkage & This is the distance between the closest members of the
two clusters &
\includegraphics[width=1.04167in,height=0.625in]{figure/chapter5/singlelinkage.png} \\
Complete Linkage & This is the distance between the members that are
farthest apart &
\includegraphics[width=1.04167in,height=0.625in]{figure/chapter5/completelinkage.png} \\
Average Linkage & This method involves looking at the distances between
all the pairs and averages all of the distances, also called UPGMA &
\includegraphics[width=1.04167in,height=0.625in]{figure/chapter5/averagelinkage.png} \\
Ward Linkage & This method minimizes the distance within the clusters
variance &
\includegraphics[width=1.04167in,height=0.625in]{figure/chapter5/wardlinkage.png} \\
\end{longtable}

Apart from the linkage methods presented in Table~\ref{tbl-linkage},
there are other procedures, such as the combined weighted algorithm
discussed by Maqbool and Babri (2004). While these other methods offer nuanced
approaches to clustering, their practical use is somewhat limited.

To illustrate the workings of HAC, consider a scenario with six elements
(referred to as `words') across two dimensions (labeled as
`documents'). This example is a slightly more complex than the
one above, but is still accessible to beginners. We have applied HAC using a single-linkage method. Notice that in the
previous example, the focus was on the clustering of documents, while now
the clustering revolves around words. We hypothesized that these words appear
with varying frequencies within the documents, meaning that in each
document, the six words appear a certain number of times. The words have
a similar characteristic if they appear with a similar frequency in the
documents that hold them.

\begin{longtable}[]{@{}lll@{}}
\caption{Table containing six words and their frequencies in the
corpus}\label{tbl-tablehac}\tabularnewline
\toprule\noalign{}
Word & Document1 & Document2 \\
\midrule\noalign{}
\endfirsthead
\toprule\noalign{}
Word & Document1 & Document2 \\
\midrule\noalign{}
\endhead
\bottomrule\noalign{}
\endlastfoot
Word1 & 40 & 53 \\
Word2 & 22 & 38 \\
Word3 & 35 & 32 \\
Word4 & 26 & 19 \\
Word5 & 8 & 41 \\
Word6 & 45 & 30 \\
\end{longtable}

The first step in starting to perform HAC is to calculate the distances
between the elements, in this case, the words. Words that exhibit similar
frequencies across documents will be positioned closer to each other in
the conceptual space. A visual representation of this arrangement, if
plotted on a canvas, would resemble the configuration shown in
Figure~\ref{fig-examplehac}.

\emph{Program code starts}

\begin{Shaded}
\begin{Highlighting}[]
\CommentTok{\# Create a data frame for six words and their frequencies in two documents}
\NormalTok{words\_data }\OtherTok{\textless{}{-}} \FunctionTok{data.frame}\NormalTok{(}
  \AttributeTok{Word =} \FunctionTok{c}\NormalTok{(}\StringTok{"Word1"}\NormalTok{, }\StringTok{"Word2"}\NormalTok{, }\StringTok{"Word3"}\NormalTok{, }\StringTok{"Word4"}\NormalTok{, }\StringTok{"Word5"}\NormalTok{, }\StringTok{"Word6"}\NormalTok{),}
  \AttributeTok{Document1 =} \FunctionTok{c}\NormalTok{(}\DecValTok{40}\NormalTok{, }\DecValTok{22}\NormalTok{, }\DecValTok{35}\NormalTok{, }\DecValTok{26}\NormalTok{, }\DecValTok{8}\NormalTok{, }\DecValTok{45}\NormalTok{),}
  \AttributeTok{Document2 =} \FunctionTok{c}\NormalTok{(}\DecValTok{53}\NormalTok{, }\DecValTok{38}\NormalTok{, }\DecValTok{32}\NormalTok{, }\DecValTok{19}\NormalTok{, }\DecValTok{41}\NormalTok{, }\DecValTok{30}\NormalTok{)}
\NormalTok{)}

\CommentTok{\# Calculate Euclidean distances}
\NormalTok{dist\_matrix }\OtherTok{\textless{}{-}} \FunctionTok{dist}\NormalTok{(words\_data[, }\SpecialCharTok{{-}}\DecValTok{1}\NormalTok{], }\AttributeTok{method =} \StringTok{"euclidean"}\NormalTok{)}
\FunctionTok{print}\NormalTok{(dist\_matrix)}

\CommentTok{\# Plotting using ggplot2}
\FunctionTok{library}\NormalTok{(ggplot2)}
\FunctionTok{ggplot}\NormalTok{(words\_data, }\FunctionTok{aes}\NormalTok{(}\AttributeTok{x =}\NormalTok{ Document1, }\AttributeTok{y =}\NormalTok{ Document2, }\AttributeTok{label =}\NormalTok{ Word)) }\SpecialCharTok{+} 
  \FunctionTok{geom\_point}\NormalTok{() }\SpecialCharTok{+} 
  \FunctionTok{geom\_text}\NormalTok{(}\AttributeTok{vjust =} \StringTok{"inward"}\NormalTok{, }\AttributeTok{hjust =} \StringTok{"inward"}\NormalTok{) }\SpecialCharTok{+}
  \FunctionTok{labs}\NormalTok{(}\AttributeTok{x =} \StringTok{"Frequency in Document 1"}\NormalTok{, }\AttributeTok{y =} \StringTok{"Frequency in Document 2"}\NormalTok{) }\SpecialCharTok{+}
  \FunctionTok{theme\_minimal}\NormalTok{()}
\end{Highlighting}
\end{Shaded}

\emph{Program code ends}

\begin{verbatim}
         1        2        3        4        5
2 23.43075                                    
3 21.58703 14.31782                           
4 36.76955 19.41649 15.81139                  
5 34.17601 14.31782 28.46050 28.42534         
6 23.53720 24.35159 10.19804 21.95450 38.60052
\end{verbatim}

\begin{figure}

\centering{

\includegraphics{5.Chapter_files/figure-pdf/fig-examplehac-1.pdf}

}

\caption{\label{fig-examplehac}Six words, two documents}

\end{figure}%

The distance between \texttt{Word1} and \texttt{Word2} can be calculated
using the Euclidean metric, \texttt{Document1} is represented on the X
axis, \texttt{Document2} on the Y axis:

\[
dist(Word1, Word2)=\sqrt{(40-22)^2+(53-35)^2}=23.43 
\] We can calculate the distance between all the pairs of
words in the same way, (Word1, Word3), (Word1, Word4), (Word1, Word5), etc., and put them
into a distance matrix, see Table~\ref{tbl-dd1}.

\begin{longtable}[]{@{}ccccccc@{}}
\caption{Distance matrix}\label{tbl-dd1}\tabularnewline
\toprule\noalign{}
& Word 1 & Word 2 & Word 3. & Word 4 & Word 5 & Word 6 \\
\midrule\noalign{}
\endfirsthead
\toprule\noalign{}
& Word 1 & Word 2 & Word 3. & Word 4 & Word 5 & Word 6 \\
\midrule\noalign{}
\endhead
\bottomrule\noalign{}
\endlastfoot
Word 1 & 0.00 & & & & & \\
Word 2 & 23.43 & 0.00 & & & & \\
Word 3 & 21.59 & 14.32 & 0.00 & & & \\
Word 4 & 36.77 & 19.42 & 15.81 & 0.00 & & \\
Word 5 & 34.18 & 14.32 & 28.46 & 28.43 & 0.00 & \\
Word 6 & 23.54 & 24.35 & \textbf{10.20} & 21.95 & 38.60 & 0.00 \\
\end{longtable}

You can see that the diagonal of the distance matrix is always zero,
because the distance between a word and itself is null. In symmetric-distance matrices the values above and below the diagonal mirror each
other, hence only one part, usually the lower, is displayed for
simplicity. For example, the distance between \texttt{Word1} and
\texttt{Word3} is the same distance as the one between
\texttt{Word3} and \texttt{Word1}. Using the Euclidean-distance formula,
we can calculate these distances and then apply single-linkage clustering,
which is also known as the minimum method. In other words, we search for the
minimum value between two elements. Initially, all the elements form
clusters on their own, resulting in six separate clusters. The minimum
value between two values is between \texttt{Word3} and \texttt{Word6},
as highlighted in Table~\ref{tbl-dd1}, leading to their combination into
the first cluster. There are two vertical lines which represent the
individual clusters, merging into a single line that represents the
new cluster. The number 1 in Figure~\ref{fig-hac1} states that this is
the first cluster, the vertical line indicates at what height on the
dendrogram the cluster was produced, in this case at \textbf{10.20}. The
height represents the distance or difference between the clusters. If
two clusters join at a low height, it suggests they are quite similar or
close to each other. Conversely, if two clusters join at a higher
height, it indicates a greater degree of dissimilarity or distance, suggesting that
the clusters are less alike.

\begin{figure}

\centering{

\includegraphics[width=0.5\textwidth,height=\textheight]{figure/chapter5/hac1.png}

}

\caption{\label{fig-hac1}Clustering Word 3 and Word 6}

\end{figure}%

After the initial clustering, we can update the distance matrix to
show the newly formed cluster, see Table~\ref{tbl-dd2}. In this
revised matrix, the distances have been recalculated to account for the clusters of
elements that have just been created. We must be consistent and apply the single-linkage criterion each time. In this updated scenario, the minimum
distance can now be seen between \texttt{Word2} and \texttt{Word5},
highlighted in bold in the table. This indicates that these two words
are the next closest pair.

\begin{longtable}[]{@{}
  >{\centering\arraybackslash}p{(\columnwidth - 10\tabcolsep) * \real{0.1667}}
  >{\centering\arraybackslash}p{(\columnwidth - 10\tabcolsep) * \real{0.1667}}
  >{\centering\arraybackslash}p{(\columnwidth - 10\tabcolsep) * \real{0.1667}}
  >{\centering\arraybackslash}p{(\columnwidth - 10\tabcolsep) * \real{0.1667}}
  >{\centering\arraybackslash}p{(\columnwidth - 10\tabcolsep) * \real{0.1667}}
  >{\centering\arraybackslash}p{(\columnwidth - 10\tabcolsep) * \real{0.1667}}@{}}
\caption{Updated distance matrix}\label{tbl-dd2}\tabularnewline
\toprule\noalign{}
\begin{minipage}[b]{\linewidth}\centering
\end{minipage} & \begin{minipage}[b]{\linewidth}\centering
Word 1
\end{minipage} & \begin{minipage}[b]{\linewidth}\centering
Word 2
\end{minipage} & \begin{minipage}[b]{\linewidth}\centering
(Word 3, Word 6)
\end{minipage} & \begin{minipage}[b]{\linewidth}\centering
Word 4
\end{minipage} & \begin{minipage}[b]{\linewidth}\centering
Word 5
\end{minipage} \\
\midrule\noalign{}
\endfirsthead
\toprule\noalign{}
\begin{minipage}[b]{\linewidth}\centering
\end{minipage} & \begin{minipage}[b]{\linewidth}\centering
Word 1
\end{minipage} & \begin{minipage}[b]{\linewidth}\centering
Word 2
\end{minipage} & \begin{minipage}[b]{\linewidth}\centering
(Word 3, Word 6)
\end{minipage} & \begin{minipage}[b]{\linewidth}\centering
Word 4
\end{minipage} & \begin{minipage}[b]{\linewidth}\centering
Word 5
\end{minipage} \\
\midrule\noalign{}
\endhead
\bottomrule\noalign{}
\endlastfoot
Word 1 & 0.00 & & & & \\
Word 2 & 23.43 & 0.00 & & & \\
(Word 3, Word 6) & (21.59, 23.54) & (14.32, 24.35) & 0.00 & & \\
Word 4 & 36.77 & 19.42 & (15.82, 21.96) & 0.00 & \\
Word 5 & 34.17 & \textbf{14.32} & (28.46, 38.60) & 28.42 & 0.00 \\
\end{longtable}

We can also focus on finding the minimum distance between the newly formed
cluster and the remaining individual words or existing clusters. To do this, we have to look for the smallest value among these recalculated
distances. In this case, the next smallest distance is observed between
\texttt{Word2} and \texttt{Word5}, indicating that these two entities
will form the next cluster. In Figure~\ref{fig-hac2} this is marked as
number 2.

\begin{figure}

\centering{

\includegraphics[width=0.5\textwidth,height=\textheight]{figure/chapter5/hac2.png}

}

\caption{\label{fig-hac2}Clustering Word 2 and Word 5}

\end{figure}%

We hope by now that the
hierarchical-clustering procedure is clear. As shown in
Figure~\ref{fig-hac5}, at the end of the process we will have one single comprehensive
final cluster, here labeled as Cluster 5. Note that the number of
clusters can be determined by cutting the dendrogram at a
specific height. To sum up, the dendrogram's horizontal axis displays
the clusters, while the vertical axis represents the sequence and scale
of cluster amalgamation, with the numerical values indicating the
proximity or similarity between the merged clusters.

\begin{figure}

\centering{

\includegraphics[width=0.5\textwidth,height=\textheight]{figure/chapter5/hac5.png}

}

\caption{\label{fig-hac5}Clustering Word 1 will al the other clusters}

\end{figure}%

Determining the optimal number of clusters in a data set is
often challenging, as there is no definitive answer. Several heuristic
methods, commonly referred to as stopping rules, may offer a degree of guidance by
generating a score or index, helping you come to a decision about the ideal number of
clusters. The choice of an appropriate number of clusters and how well
data are clustered is subjective and influenced by various factors,
including which distance metric is used.

Some of the prominent stopping rules or indices dealing with clustering
validation include the Bayesian information criterion (BIC), Akaike
information criterion (AIC), silhouette score, Davies-Bouldin index,
homogeneity score, completeness score, V-measure, the Calinski-Harabasz
index, Duda-Hartigan's rule, among others. It is crucial to note that while
these mathematical methods can deliver an optimal result, the clusters
formed may not always align with intuitive or meaningful groupings in
practical scenarios.

In applying HAC to our dataset, we have introduced variation. Instead of using the
conventional Euclidean distance, we opted for the Canberra distance
(Kocher and Savoy 2017), and instead of single linkage, we employ Ward's
method (Strauss and Maltitz 2017). As we have already mentioned, we can
also employ the cosine distance (Maimon and Rokach 2010). For assessing
the optimal number of clusters, we used the Hartigan index
(Purwitasari et al. 2011) which calculates the maximum difference
between the levels of the index's hierarchy.

\emph{Program code starts}

\begin{Shaded}
\begin{Highlighting}[]
\CommentTok{\# Loading libraries}
\FunctionTok{library}\NormalTok{(factoextra)}
\FunctionTok{library}\NormalTok{(NbClust)}
\FunctionTok{library}\NormalTok{(proxy)}

\CommentTok{\# Convert \textquotesingle{}docs\textquotesingle{} to matrix}
\NormalTok{m }\OtherTok{\textless{}{-}} \FunctionTok{as.matrix}\NormalTok{(docs)}

\CommentTok{\# Canberra distance matrix}
\NormalTok{d }\OtherTok{\textless{}{-}} \FunctionTok{dist}\NormalTok{(m, }\AttributeTok{method =} \StringTok{"canberra"}\NormalTok{)}

\CommentTok{\# Ward\textquotesingle{}s method clustering}
\NormalTok{hc }\OtherTok{\textless{}{-}} \FunctionTok{hclust}\NormalTok{(d, }\AttributeTok{method =} \StringTok{"ward.D"}\NormalTok{)}

\CommentTok{\# Optimal cluster count}
\NormalTok{nb }\OtherTok{\textless{}{-}} \FunctionTok{NbClust}\NormalTok{(m, }\AttributeTok{distance =} \StringTok{"canberra"}\NormalTok{, }\AttributeTok{min.nc =} \DecValTok{2}\NormalTok{, }\AttributeTok{max.nc =} \DecValTok{5}\NormalTok{, }
              \AttributeTok{method =} \StringTok{"ward.D"}\NormalTok{, }\AttributeTok{index =} \StringTok{"all"}\NormalTok{)}
\end{Highlighting}
\end{Shaded}

*** : The Hubert index is a graphical method of determining the number
of clusters. In the plot of a Hubert index, we are looking for a significant knee
that corresponds to a significant increase in the value of the measure,
i.e., the significant peak in Hubert-index second differences plot.

*** : The D index is a graphical method of determining the number of
clusters. In the plot of a D index, we are looking for a significant knee (the
significant peak in the D index second differences plot) that corresponds to
a significant increase in the value of the measure.

\begin{center}\rule{0.5\linewidth}{0.5pt}\end{center}

\begin{itemize}
\item
  Of the indices:\\
\item
  6 proposed 2 as the best number of clusters
\item
  9 proposed 3 as the best number of clusters
\item
  1 proposed 4 as the best number of clusters
\item
  6 proposed 5 as the best number of clusters

\begin{verbatim}
             ***** Conclusion *****                            
\end{verbatim}
\item
  According to the majority, the best number of clusters is 3
\end{itemize}

\begin{center}\rule{0.5\linewidth}{0.5pt}\end{center}

\begin{Shaded}
\begin{Highlighting}[]
\NormalTok{optimal\_clusters }\OtherTok{\textless{}{-}}\NormalTok{ nb}\SpecialCharTok{$}\NormalTok{Best.nc[}\DecValTok{1}\NormalTok{]}

\CommentTok{\# Clustering result}
\NormalTok{hc\_result }\OtherTok{\textless{}{-}} \FunctionTok{hcut}\NormalTok{(d, }\AttributeTok{k =}\NormalTok{ optimal\_clusters, }\AttributeTok{hc\_method =} \StringTok{"ward.D"}\NormalTok{)}
\end{Highlighting}
\end{Shaded}

\emph{Program code ends}

To visualize the hierarchical clustering results as a dendrogram, see
Figure~\ref{fig-dendrogram},. To do this we used the function
\texttt{fviz\_dend()} from the factoextra package (Kassambara and Mundt 2020). It uses \texttt{ggplot2} to create aesthetically informative dendrograms. For the calculation of
the distances between the clustered items, we used the \texttt{dist()}
function, with the cosine distance measure from the
\texttt{proxy} package (Meyer and Buchta 2021).

\emph{Program code starts}

\begin{Shaded}
\begin{Highlighting}[]
\CommentTok{\# Dendrogram visualization}
\FunctionTok{fviz\_dend}\NormalTok{(hc\_result, }\AttributeTok{rect =} \ConstantTok{TRUE}\NormalTok{, }\AttributeTok{rect\_border =} \StringTok{"jco"}\NormalTok{, }\AttributeTok{rect\_lty =} \DecValTok{2}\NormalTok{, }\AttributeTok{cex =} \FloatTok{0.5}\NormalTok{)}
\end{Highlighting}
\end{Shaded}

\begin{figure}[H]

\centering{

\includegraphics{5.Chapter_files/figure-pdf/fig-dendrogram-1.pdf}

}

\caption{\label{fig-dendrogram}Practical example}

\end{figure}%

\emph{Program code ends}

We can also visualize the cluster plot too, see
Figure~\ref{fig-clusterplot}.

\emph{Program code starts}

\begin{Shaded}
\begin{Highlighting}[]
\CommentTok{\# Cluster visualization}
\FunctionTok{fviz\_cluster}\NormalTok{(hc\_result, }\AttributeTok{data=}\NormalTok{m, }\AttributeTok{ellipse.type =} \StringTok{"norm"}\NormalTok{)}
\end{Highlighting}
\end{Shaded}

\begin{figure}[H]

\centering{

\includegraphics{5.Chapter_files/figure-pdf/fig-clusterplot-1.pdf}

}

\caption{\label{fig-clusterplot}Practical example}

\end{figure}%

\emph{Program code ends}

In this analysis, we used a weighted version of the term-document
matrix, see Chapter~\ref{sec-basic-text-analysis-and-statistics}. Unlike
the regular TDM, which employs plain word frequencies, the weighted
version contains a numerical statistic and considers the significance of
words across the entire corpus.

The removal of most of the words was done to manage the complexity and
ensure clarity of the visualization. Try doing this without the removal
of the words (inadvisable) to see the chaos that ensues.

The clustering process grouped words that had a similar
importance across the documents. Yes, \emph{kindle fire} is one such case,
and as you can see the words \emph{kindle} and \textit{fire} frequently occur together
in the context of this dataset. It is the experience and the knowledge
of the analyst that make sense of the clusters and the underlying semantic
relationships and patterns.

The \texttt{NbClust()} function yielded 3 as the best number of clusters
using the Hubert index. Ideally, we would apply all the indexes to the
function and choose the mode (the most frequent number). Unfortunately
not all of them are suitable for hierarchical clustering.
Nevertheless, for all those that can be treated in this way, the mode is 3.

The choice of the best linking criterion or distance metric in
hierarchical clustering is subject to debate, as highlighted in sources
such as Gower and Legendre (1986). There is a consensus that the choice
largely depends on the specific circumstances of the data and the
analyst's experience. As we have already mentioned, Euclidean
distance might not always be appropriate as it tends to group documents
based on length rather than semantic content. In such cases, cosine
similarity, which focuses on the orientation rather than the magnitude
of vectors, can offer a more nuanced understanding of the textual data.
Here is an example which demonstrates how these changes can shape the final
result.

\begin{Shaded}
\begin{Highlighting}[]
\FunctionTok{library}\NormalTok{(proxy)  }\CommentTok{\# for cosine similarity calculation}
\FunctionTok{library}\NormalTok{(cluster)  }\CommentTok{\# for hierarchical clustering}

\CommentTok{\# Compute the similarity matrix}
\NormalTok{similarity\_matrix }\OtherTok{\textless{}{-}} \FunctionTok{simil}\NormalTok{(m, }\AttributeTok{method =} \StringTok{"cosine"}\NormalTok{)}

\CommentTok{\# Converting similarity to distance}
\NormalTok{distance\_matrix }\OtherTok{\textless{}{-}} \DecValTok{1} \SpecialCharTok{{-}}\NormalTok{ similarity\_matrix}

\CommentTok{\# Performing hierarchical clustering}
\NormalTok{hc }\OtherTok{\textless{}{-}} \FunctionTok{hclust}\NormalTok{(}\FunctionTok{as.dist}\NormalTok{(distance\_matrix), }\AttributeTok{method =} \StringTok{"ward.D"}\NormalTok{)  }\CommentTok{\# or any other method}

\CommentTok{\# Plotting the dendrogram}
\FunctionTok{fviz\_dend}\NormalTok{(hc, }\AttributeTok{cex =} \FloatTok{0.5}\NormalTok{, }\AttributeTok{rect =} \ConstantTok{TRUE}\NormalTok{, }\AttributeTok{k\_colors =} \StringTok{"jco"}\NormalTok{, }\AttributeTok{rect\_border =} \StringTok{"jco"}\NormalTok{, }\AttributeTok{rect\_lty =} \DecValTok{2}\NormalTok{)}
\end{Highlighting}
\end{Shaded}

\includegraphics{5.Chapter_files/figure-pdf/unnamed-chunk-1-1.pdf}

\section{K-means clustering}\label{k-means-clustering}

K-means clustering is a partitioning method that segments a data set into
a predetermined number of clusters, denoted by 'k'. The data set is divided into
partitions, where each data point belongs to only one group. The
objective is to partition the data set into k clusters to maximize the similarity of the data within each cluster, and minimize that
between clusters. To do this, it plays with the concepts of between and within.
The process begins by specifying `k,' although finding the optimal `k'
can be challenging, and methods like the elbow or silhouette technique
are often employed (see, Saputra, Saputra, and Oswari 2020). The algorithm
randomly selects the intial `k' centroids (cluster centers), either from
the data points or random locations. Each point is then assigned to its
nearest centroid based on Euclidean distance (or other distance
metrics). The centroids are recalculated as the mean of all the points within their
cluster,s and points are reassigned to the closest new centroids. This
process repeats until points no longer switch clusters or a maximum
number of iterations has been reached (typically 10 in R). The goal is to
minimize within-cluster variance, making k-means suitable for spherical
clusters; it may face challenges when confronted by clusters of varying densities or
non-spherical shapes.

Let us apply the k-means algorithm to our data set consisting of
six words and two documents. Because k-means clustering involves a degree of
randomness in its initial assignments, this can lead to
different outcomes from separate runs. The randomness is managed in R
using the \texttt{kmeans()} function's \texttt{nstart} parameter, which
carries out multiple initial configurations (25 in our case) to find the
best starting point. This approach helps mitigate the algorithm's
tendency to converge to local optima. To make the algorithm's process
easier to understand, we will use the \texttt{animation} library in R to
create a visual clustering sequence.

For our example, we will place the data into a table (using the \texttt{dt}
variable) and choose to segment it into two clusters, showing how
k-means clustering partitions the data without the hierarchical structure seen in
other methods such as HAC. The version printed here is not that
spectacular, you see only the final result, but once you have run the
code on you machine, you will appreciate and understand the power of
animation as an educational tool.

\emph{Program code starts}

\begin{Shaded}
\begin{Highlighting}[]
\FunctionTok{library}\NormalTok{(animation)}

\CommentTok{\# Example, six words with their frequencies}
\NormalTok{Rowname }\OtherTok{\textless{}{-}} \FunctionTok{c}\NormalTok{(}\StringTok{"Word1"}\NormalTok{, }\StringTok{"Word2"}\NormalTok{, }\StringTok{"Word3"}\NormalTok{, }\StringTok{"Word4"}\NormalTok{, }\StringTok{"Word5"}\NormalTok{, }\StringTok{"Word6"}\NormalTok{)}
\NormalTok{Document1 }\OtherTok{\textless{}{-}} \FunctionTok{c}\NormalTok{(}\DecValTok{40}\NormalTok{, }\DecValTok{22}\NormalTok{, }\DecValTok{35}\NormalTok{, }\DecValTok{26}\NormalTok{, }\DecValTok{8}\NormalTok{, }\DecValTok{45}\NormalTok{)}
\NormalTok{Document2 }\OtherTok{\textless{}{-}} \FunctionTok{c}\NormalTok{(}\DecValTok{53}\NormalTok{, }\DecValTok{38}\NormalTok{, }\DecValTok{32}\NormalTok{, }\DecValTok{19}\NormalTok{, }\DecValTok{41}\NormalTok{, }\DecValTok{30}\NormalTok{)}

\CommentTok{\# Create the dataset}
\NormalTok{dt }\OtherTok{\textless{}{-}} \FunctionTok{cbind}\NormalTok{(Document1, Document2) }\SpecialCharTok{\%\textgreater{}\%}
\NormalTok{  as.data.frame }\SpecialCharTok{\%\textgreater{}\%}
  \StringTok{\textasciigrave{}}\AttributeTok{rownames\textless{}{-}}\StringTok{\textasciigrave{}}\NormalTok{(Rowname)}

\CommentTok{\# Set a random seed for reproducibility}
\FunctionTok{set.seed}\NormalTok{(}\DecValTok{123}\NormalTok{)}

\CommentTok{\# Perform animated k{-}means with 2 clusters}
\FunctionTok{kmeans.ani}\NormalTok{(dt, }\DecValTok{2}\NormalTok{)}
\end{Highlighting}
\end{Shaded}

\includegraphics{5.Chapter_files/figure-pdf/kmeansani-1.pdf}

\includegraphics{5.Chapter_files/figure-pdf/kmeansani-2.pdf}

\includegraphics{5.Chapter_files/figure-pdf/kmeansani-3.pdf}

\includegraphics{5.Chapter_files/figure-pdf/kmeansani-4.pdf}

\emph{Program code ends}

Several limitations need to be taken into account for k-means clustering. One is
that the algorithm may perform poorly when cluster sizes and densities
vary significantly. Another one is that the presence of outliers can
skew the clustering results, as k-means clustering is sensitive to extreme
values. Additionally, this technique assumes clusters are roughly spherical
in shape, meaning it struggles with non-spherical (e.g., elongated or
oblong) clusters. In a two-dimensional space, k-means is most effective
when data points within a cluster are equidistant from the cluster
center, forming a circular pattern. However, this occurs rarely in real-world data, as illustrated by the oblong shape of cluster 1
in Figure~\ref{fig-oblongandnot}, which is not ideal for for this form of clustering.

\begin{figure}

\centering{

\includegraphics[width=0.99\textwidth,height=\textheight]{figure/chapter5/oblongandnot.png}

}

\caption{\label{fig-oblongandnot}k-means clustering situations}

\end{figure}%

In our example, we analyzed the reviews data set
using the k-means clustering algorithm, employing the \texttt{kmeans()}
function in R. However, we let \texttt{NbClust()} find the optima number of
clusters. The \texttt{nstart} parameter sets the number of initial random centroid configurations to try, with the algorithm selecting the
configuration that yields the most stable arrangement of clusters. Empirical
observations suggest that a higher value often benefits smaller data sets
by enhancing the stability of the clusters. We use the \texttt{fviz\_cluster()}
function to visualize the clustering results, which offers several
helpful formatting options. The \texttt{repel} argument helps to prevent
label overlap, enhancing readability, while \texttt{show.clust.cent}
plots the centroid of each cluster, providing a clear visual
representation of the clustering. After plotting, it is possible to examine the cluster assignment of individual words through the cluster
attribute in the \texttt{km} object (accessible via
\texttt{km\$cluster}). This analysis allows us to understand how words
from the reviews can be grouped into distinct clusters.

\begin{Shaded}
\begin{Highlighting}[]
\FunctionTok{library}\NormalTok{(RColorBrewer)}
\FunctionTok{set.seed}\NormalTok{(}\DecValTok{123}\NormalTok{)}


\NormalTok{tdm\_matrix }\OtherTok{\textless{}{-}} \FunctionTok{as.matrix}\NormalTok{(docs)}

\CommentTok{\# Calculate the optimal number of clusters using NbClust}
\NormalTok{nckm }\OtherTok{\textless{}{-}} \FunctionTok{NbClust}\NormalTok{(tdm\_matrix, }\AttributeTok{distance =} \StringTok{"canberra"}\NormalTok{, }\AttributeTok{min.nc =} \DecValTok{2}\NormalTok{, }\AttributeTok{max.nc =} \DecValTok{10}\NormalTok{, }
                \AttributeTok{method =} \StringTok{"kmeans"}\NormalTok{, }\AttributeTok{index =} \StringTok{"ch"}\NormalTok{)}
\NormalTok{nbc }\OtherTok{\textless{}{-}} \FunctionTok{as.vector}\NormalTok{(nckm}\SpecialCharTok{$}\NormalTok{Best.nc[}\DecValTok{1}\NormalTok{])  }\CommentTok{\# Extracting the optimal cluster number}

\CommentTok{\# Perform k{-}means clustering with the optimal number of clusters}
\NormalTok{km }\OtherTok{\textless{}{-}} \FunctionTok{kmeans}\NormalTok{(tdm\_matrix, }\AttributeTok{centers =}\NormalTok{ nbc, }\AttributeTok{nstart =} \DecValTok{30}\NormalTok{)}

\CommentTok{\# Generate a color palette based on the number of clusters}
\NormalTok{palette }\OtherTok{\textless{}{-}} \FunctionTok{brewer.pal}\NormalTok{(nbc, }\StringTok{"Set1"}\NormalTok{)}

\CommentTok{\# Visualize the clusters using the same data matrix used in kmeans}
\FunctionTok{fviz\_cluster}\NormalTok{(}\FunctionTok{list}\NormalTok{(}\AttributeTok{data =}\NormalTok{ tdm\_matrix, }\AttributeTok{cluster =}\NormalTok{ km}\SpecialCharTok{$}\NormalTok{cluster),}
             \AttributeTok{palette =}\NormalTok{ palette,}
             \AttributeTok{main =} \StringTok{"Partitioning Clustering Plot"}\NormalTok{,}
             \AttributeTok{ellipse.type =} \StringTok{"norm"}\NormalTok{,}
             \AttributeTok{repel =} \ConstantTok{TRUE}\NormalTok{,}
             \AttributeTok{show.clust.cent =} \ConstantTok{TRUE}
\NormalTok{)}
\end{Highlighting}
\end{Shaded}

\includegraphics{5.Chapter_files/figure-pdf/kclustering-1.pdf}

\begin{Shaded}
\begin{Highlighting}[]
\CommentTok{\# Print cluster assignments}
\NormalTok{km}\SpecialCharTok{$}\NormalTok{cluster}
\end{Highlighting}
\end{Shaded}

\begin{verbatim}
     able    access      also    amazon   android  anything       app      apps 
        1         1         1         1         1         1         1         1 
     back   battery    better    bought       buy       can      cant  computer 
        1         1         1         2         1         2         1         1 
      day    device    doesnt      dont      easy     email      even      fire 
        1         1         1         1         1         1         1         2 
    first       get      good       got     great  internet      ipad       ive 
        1         1         2         1         2         2         1         1 
     just      keep    kindle      like    little      long   looking      love 
        1         1         2         1         1         1         1         2 
     many      much      need     never      nice       now       one     price 
        1         1         1         1         1         1         2         1 
  product purchased      read    really  received recommend   reviews    screen 
        1         1         2         1         1         1         2         1 
      see     since something     still     store   support    tablet      take 
        1         1         1         1         1         1         1         1 
    thing     think      time       use      used     using      want    wanted 
        1         1         1         1         1         1         1         1 
      way      well      wifi      will      work     works     youre 
        1         1         1         1         2         2         1 
\end{verbatim}

We determined our data set's optimal number of clusters using the k-means
method and by applying the \texttt{ch} index. This resulted in an
optimal cluster count of 2. This process enables us to visually inspect
and understand the distribution of words across the clusters.

The k-means algorithm has several variations tailored to different characteristics of the data. One such variation is k-medians, where the cluster's
center is determined by the median rather than the mean, offering a
robust defense against outliers (Arora, Raghavan, and Rao 1998). Another
variation, k-medoids, designates a real data point as the centre of the cluster, usually the most centrally located point in the cluster,
enhancing the clarity of the interpretation of the clusters (Rdusseeun and Kaufman
1987). Spherical k-means clustering is particularly useful for high-dimensional
data. In this approach, data points are considered to be vectors on a unit
sphere, and the cosine similarity is used as the distance metric, making
it suitable for text and other high-dimensional data (Dhillon and Modha
2001).

\subsection{Some considerations regarding
k-clustering}\label{considerations-on-k-clustering}

K-means clustering is a widely used algorithm due to its simplicity and
computational efficiency, making it suitable for large data sets with
multiple variables. The algorithm's core requirement is that the
number of clusters (k) is specified, and it adapts to new partitions by
simply changing this value. It performs efficiently when clusters are
spherical or circular in a 2D plane, often resulting in tightly grouped
clusters allowing for clear interpretation when compared to hierarchical clustering.

However, selecting the optimal number of clusters can be challenging,
and results may vary between runs due to the algorithm's sensitivity to the
initial conditions. This inconsistency is a notable limitation.
Moreover, it typically assumes the clusters are of approximately equal
sizes, which may not align with the underlying distribution of the data.
Another limitation is its exclusive focus on numerical data, which
relies on calculating the Euclidean (or other types of) distance between
points.

\section{Gaussian mixtures clustering}\label{gaussian-mixtures-clusterig}

Gaussian mixture clustering (GMC), often used in conjunction with
the expectation-maximization (EM) algorithm, is a sophisticated method
for identifying underlying patterns in data (Bouman et al. 1997). This
approach assumes that data points are generated using a combination of
several Gaussian (or normal) distributions, each representing a distinct
cluster. The term `mixture' in GMC reflects this concept of multiple,
potentially overlapping, Gaussian distributions within the data set.
Unlike the `hard' approach of k-means clustering, where each data point
is assigned to one cluster, GMC adopts a `soft clustering'
technique. Each data point is associated with probabilities denoting its
likelihood of belonging to different clusters, allowing for more nuanced
and probabilistic classifications. In GMC, each Gaussian distribution
within the model is defined by its own mean and covariance, shaping the
cluster's appearance in the data space. The EM algorithm plays an
essential role in GMC by iteratively estimating and refining these
parameters, aiming to maximize the fit of the model to the
data.

This model is handy in scenarios where data does not conform to the
spherical cluster assumed of k-means clustering, as shown by cluster 1 of
Figure~\ref{fig-oblongandnot}. For instance, GMC can handle
elongated or irregularly shaped clusters effectively. In R, the Gaussian mixture
model can be implemented using the \texttt{GMM()} function from the
\texttt{ClusterR} package (Mouselimis 2023). The process typically
involves converting data into a suitable format (such as a data frame) and
specifying the number of clusters. The \texttt{predict\_GMM}
function assigns each data point to clusters based on the probabilities
calculated by \texttt{GMM()}. It is advisable to explore the functions'
parameters and settings to understand their application. We will now explore
this type of clustering using 3 clusters.

\begin{Shaded}
\begin{Highlighting}[]
\FunctionTok{library}\NormalTok{(ClusterR)}

\CommentTok{\# Convert \textquotesingle{}docs\textquotesingle{} to distance matrix and then to a data frame}
\NormalTok{dist\_matrix\_df }\OtherTok{\textless{}{-}} \FunctionTok{as.matrix}\NormalTok{(docs) }\SpecialCharTok{\%\textgreater{}\%}
                  \FunctionTok{dist}\NormalTok{() }\SpecialCharTok{\%\textgreater{}\%}
                  \FunctionTok{as.matrix}\NormalTok{() }\SpecialCharTok{\%\textgreater{}\%}
                  \FunctionTok{as.data.frame}\NormalTok{()}

\CommentTok{\# Set the number of clusters}
\NormalTok{nbc }\OtherTok{\textless{}{-}} \DecValTok{3}

\CommentTok{\# Create Gaussian mixture model (use Mahalanobis distance)}
\NormalTok{gmc\_model }\OtherTok{\textless{}{-}} \FunctionTok{GMM}\NormalTok{(dist\_matrix\_df, }\AttributeTok{gaussian\_comps =}\NormalTok{ nbc, }
                 \AttributeTok{dist\_mode =} \StringTok{"maha\_dist"}\NormalTok{, }\AttributeTok{seed\_mode =} \StringTok{"random\_subset"}\NormalTok{,}
                 \AttributeTok{km\_iter =} \DecValTok{10}\NormalTok{, }\AttributeTok{em\_iter =} \DecValTok{5}\NormalTok{, }\AttributeTok{verbose =} \ConstantTok{FALSE}\NormalTok{,}
                 \AttributeTok{var\_floor =} \FloatTok{1e{-}10}\NormalTok{, }\AttributeTok{seed =} \DecValTok{1}\NormalTok{)}

\CommentTok{\# Predict cluster labels}
\NormalTok{gmc\_labels }\OtherTok{\textless{}{-}} \FunctionTok{predict\_GMM}\NormalTok{(dist\_matrix\_df, }\AttributeTok{CENTROIDS =}\NormalTok{ gmc\_model}\SpecialCharTok{$}\NormalTok{centroids, }
                          \AttributeTok{COVARIANCE =}\NormalTok{ gmc\_model}\SpecialCharTok{$}\NormalTok{covariance\_matrices, }
                          \AttributeTok{WEIGHTS =}\NormalTok{ gmc\_model}\SpecialCharTok{$}\NormalTok{weights)}\SpecialCharTok{$}\NormalTok{cluster\_labels}

\CommentTok{\# Visualize the clusters}
\FunctionTok{fviz\_cluster}\NormalTok{(}\FunctionTok{list}\NormalTok{(}\AttributeTok{data =}\NormalTok{ dist\_matrix\_df, }\AttributeTok{cluster =}\NormalTok{ gmc\_labels), }
             \AttributeTok{ellipse.type =} \StringTok{"norm"}\NormalTok{, }\AttributeTok{repel =} \ConstantTok{TRUE}\NormalTok{)}
\end{Highlighting}
\end{Shaded}

\includegraphics{5.Chapter_files/figure-pdf/gaussian-1.pdf}

\begin{Shaded}
\begin{Highlighting}[]
\CommentTok{\# Create a data frame of words and their cluster labels}
\NormalTok{df }\OtherTok{\textless{}{-}} \FunctionTok{data.frame}\NormalTok{(}\AttributeTok{words =} \FunctionTok{rownames}\NormalTok{(dist\_matrix\_df), }\AttributeTok{cluster =}\NormalTok{ gmc\_labels)}

\CommentTok{\# Display some results}
\FunctionTok{print}\NormalTok{(}\FunctionTok{head}\NormalTok{(df))}
\end{Highlighting}
\end{Shaded}

\begin{verbatim}
     words cluster
1     able       3
2   access       3
3     also       3
4   amazon       2
5  android       1
6 anything       3
\end{verbatim}

Here we can see better clustering compared to that produced by k-means clustering. It is also possible to check which clusters the words belong to.

\subsection{Considerations regarding GMM}\label{considerations-on-gmm}

GMM introduces significant degree of flexibility compared to k-means clustering, and particularly
in accommodating clusters of different shapes, orientations, and
variances. This is partly due to GMM's ability to model each cluster
with its own covariance structure. Unlike the hard clustering of the
k-means approach, where each object is assigned to a single cluster, GMM allows
for probabilistic assignments to clusters. This soft-clustering approach is
particularly advantageous in text analytics where words often relate to
multiple topics. In GMM, centroid updates consider both the mean and
variance, capturing more complex data structures than the mean-only
updates of k-means clustering. However, GMM assumes that you have Gaussian distributed data and
may have difficulties working with highly correlated datasets. It also requires
careful selection of the number of clusters and can be computationally
demanding, particularly for large datasets. Despite these limitations,
GMM is a powerful algorithm that performs well, especially where data
exhibits complex, overlapping clustering patterns.

\section{Density-based clustering}\label{density-based-clustering}

DBSCAN (density-based spatial clustering of applications with noise) is
ideal for datasets where clusters are not necessarily spherical,
offering, like GMC, a flexible alternative to methods such as k-means clustering.
It operates using two parameters: a radius \(\epsilon\) and a minimum
number of points (minPts). DBSCAN forms clusters from dense regions of
data points, delineated by the \(\epsilon\) radius, and identifies
outliers in low-density areas. Unlike some algorithms, it does not
require you to specify the number of clusters, as these
emerge naturally from the data based on its density. It is
important to know that if a data point falls within the \(\epsilon\)
radius of multiple clusters, it is typically assigned to the first
cluster formed. Choosing the right \(\epsilon\) can be challenging, as
it significantly affects the resulting clusters, but there are methods
to find both \(\epsilon\) and minPts (Ozkok and Celik 2017).

HDBSCAN (hierarchical density-based spatial clustering of applications
with noise) extends DBSCAN's capabilities by transforming its flat
clustering into a hierarchical structure. This method extracts clusters
based on their stability over varying levels of density, with no need for
the user to specify \(\epsilon\). It is particularly good at handling
datasets with clusters of different densities, automatically determining
the optimal scale for \(\epsilon\) and revealing consistent clusters
over a range of densities.

We will use the HDBSCAN algorithm, available in the \texttt{dbscan} package, to explore density-based clustering. This approach requires setting
hyper-parameters, which often depend on the data set and objectives of the
clustering. HDBSCAN's ability to identify stable clusters
without predefined parameters like \(\epsilon\) makes it a robust and
user-friendly choice when confronted with complex datasets.

\begin{Shaded}
\begin{Highlighting}[]
\FunctionTok{library}\NormalTok{(dbscan)}

\CommentTok{\# calculate the hdbscan cluster}
\NormalTok{clustering.dbscan }\OtherTok{\textless{}{-}} \FunctionTok{hdbscan}\NormalTok{(docs, }\AttributeTok{minPts =} \DecValTok{2}\NormalTok{)}
\NormalTok{clustering.dbscan}\SpecialCharTok{$}\NormalTok{cluster}
\end{Highlighting}
\end{Shaded}

\begin{verbatim}
 [1] 2 2 2 2 0 2 2 2 2 2 2 0 2 2 2 2 2 0 2 2 2 2 2 0 2 2 0 2 2 2 0 2 2 2 0 2 2 2
[39] 2 0 2 2 2 2 1 2 0 2 0 1 0 2 2 1 0 2 2 2 2 2 2 2 2 2 2 2 2 2 2 2 2 2 2 2 2 2
[77] 0 2 2
\end{verbatim}

\begin{Shaded}
\begin{Highlighting}[]
\CommentTok{\# visualize the clusters created with hdbscan}
\FunctionTok{fviz\_cluster}\NormalTok{(}\FunctionTok{list}\NormalTok{(}\AttributeTok{data =} \FunctionTok{as.matrix}\NormalTok{(docs), }\AttributeTok{cluster =}\NormalTok{ clustering.dbscan}\SpecialCharTok{$}\NormalTok{cluster), }
             \AttributeTok{ellipse.type =} \StringTok{"norm"}\NormalTok{,}
             \AttributeTok{repel =} \ConstantTok{TRUE}\NormalTok{)}
\end{Highlighting}
\end{Shaded}

\includegraphics{5.Chapter_files/figure-pdf/dbclustering-1.pdf}

\begin{Shaded}
\begin{Highlighting}[]
\FunctionTok{library}\NormalTok{(dbscan)}
\FunctionTok{library}\NormalTok{(factoextra)}

\CommentTok{\# Perform HDBSCAN clustering}
\NormalTok{clustering\_hdbscan }\OtherTok{\textless{}{-}} \FunctionTok{hdbscan}\NormalTok{(}\FunctionTok{as.matrix}\NormalTok{(docs), }\AttributeTok{minPts =} \DecValTok{2}\NormalTok{)}

\CommentTok{\# Check clustering results}
\NormalTok{clustering\_hdbscan}
\end{Highlighting}
\end{Shaded}

\begin{verbatim}
HDBSCAN clustering for 79 objects.
Parameters: minPts = 2
The clustering contains 2 cluster(s) and 13 noise points.

 0  1  2 
13  3 63 

Available fields: cluster, minPts, coredist, cluster_scores,
                  membership_prob, outlier_scores, hc
\end{verbatim}

\begin{Shaded}
\begin{Highlighting}[]
\CommentTok{\# Create a list structure that is compatible with fviz\_cluster}
\NormalTok{cluster\_data }\OtherTok{\textless{}{-}} \FunctionTok{list}\NormalTok{(}\AttributeTok{data =} \FunctionTok{as.matrix}\NormalTok{(docs), }\AttributeTok{cluster =}\NormalTok{ clustering\_hdbscan}\SpecialCharTok{$}\NormalTok{cluster)}

\CommentTok{\# Visualize the clusters using fviz\_cluster}
\FunctionTok{fviz\_cluster}\NormalTok{(cluster\_data, }\AttributeTok{ellipse.type =} \StringTok{"norm"}\NormalTok{, }\AttributeTok{repel =} \ConstantTok{TRUE}\NormalTok{)}
\end{Highlighting}
\end{Shaded}

\includegraphics{5.Chapter_files/figure-pdf/dbclustering-2.pdf}

We can see how it has created a big cluster and two smaller ones. With
a bigger data set the clustering would yield better results.

\subsection{Considerations regarding (H)DBSCAN clustering}\label{hdbscan-clustering-considerations.}

DBSCAN and its variant HDBSCAN separate clusters into those of high- and low-density, but face difficulties when the densities are similar. Unlike in
other clustering algorithms, the clusters can take any shape, outliers are dealt with well, and there is no need to set the
number of clusters manually. However, you must tune the \(\epsilon\) and the
minimum points parameters. The order of the data has an influence
on the final result and you need sufficient data to separate dense
areas, otherwise you obtain poor results.

As a practical example, it is possible to apply DBSCAN
to the data set and find clusters based on the reviews. Using the clusters we can see similarities between reviews, for example,
if one customer writes about the screen, memory, and connection of a
tablet, and another one about the screen and memory, then it is useful to
enquire about the connection between the two. In a commercial setup, if one customer
talks about products 1, 2 and 3, and another one only about the first
two, a system of recommendation might propose the third product to the second
customer.

\section{Affinity propagation
clustering}\label{affinity-propagation-clustering}

Affinity propagation clustering, as defined by Frey and Dueck (2007 and 2009), introduces the concept of
``exemplars''. These are the foundational elements, or the core nodes, of
a data set around which clustering is performed. The idea is that all the elements are potential exemplars and a message
system between data points identifies a subset of exemplars
and their clusters iteratively.

This is carried out by the \texttt{apcluster} package
(Bodenhofer, Kothmeier, and Hochreiter 2011). Unlike in the other
clustering algorithms, the data can be of any shape and of variable
density. Importantly, the basic version of apc will
calculate the optimal number of clusters without any specifications being given.
Nevertheless, the \texttt{apclusterK()} function in the
\texttt{apcluster} package provides a mechanism to specify the cluster
count. In this, the first parameter, \texttt{s},
calculates the similarity function using the negative squared distances
by default, although alternative similarity measures, such as the Spearman
correlation, can be employed to capture the relationships within certain
datasets.

The second parameter, \texttt{x}, represents the distance matrix we
derived from on our previous analysis using k-means clustering. Despite the
possibility of setting a predetermined number of clusters via \(K\), the
adaptive nature of the apcluster() function's algorithm can ascertain
the optimal cluster count.

\begin{Shaded}
\begin{Highlighting}[]
\CommentTok{\# affinity propagation}
\FunctionTok{library}\NormalTok{(apcluster)}
\FunctionTok{library}\NormalTok{(reshape2)}


\CommentTok{\# negDistMat computes the similarity matrix }
\NormalTok{similarity\_matrix }\OtherTok{\textless{}{-}} \FunctionTok{negDistMat}\NormalTok{(docs, }\AttributeTok{r=}\DecValTok{2}\NormalTok{) }

\CommentTok{\# Apply clustering}
\NormalTok{apck }\OtherTok{\textless{}{-}} \FunctionTok{apclusterK}\NormalTok{(}\AttributeTok{s=}\NormalTok{similarity\_matrix, }\AttributeTok{K=}\DecValTok{3}\NormalTok{)}
\end{Highlighting}
\end{Shaded}

\begin{verbatim}
Trying p = -0.006489952 
   Number of clusters: 75 
Trying p = -0.01611412 
   Number of clusters: 53 
Trying p = -0.1123558 
   Number of clusters: 4 
Trying p = -0.5400968 (bisection step no. 1 )
   Number of clusters: 1 
Trying p = -0.2727587 (bisection step no. 2 )
   Number of clusters: 1 
Trying p = -0.1390897 (bisection step no. 3 )
   Number of clusters: 3 

Number of clusters: 3 for p = -0.1390897 
\end{verbatim}

\begin{Shaded}
\begin{Highlighting}[]
\CommentTok{\# Convert data from a cluster type into a data frame}
\NormalTok{apckmelt }\OtherTok{\textless{}{-}} \FunctionTok{melt}\NormalTok{(apck}\SpecialCharTok{@}\NormalTok{clusters)}

\CommentTok{\# Visualize the data}
\FunctionTok{fviz\_cluster}\NormalTok{(}\FunctionTok{list}\NormalTok{(}\AttributeTok{data =} \FunctionTok{as.matrix}\NormalTok{(docs)[apckmelt}\SpecialCharTok{$}\NormalTok{value,], }\AttributeTok{cluster =}\NormalTok{ apckmelt}\SpecialCharTok{$}\NormalTok{L1), }
             \AttributeTok{ellipse.type =} \StringTok{"norm"}\NormalTok{,}
             \AttributeTok{repel =} \ConstantTok{TRUE}\NormalTok{)}
\end{Highlighting}
\end{Shaded}

\includegraphics{5.Chapter_files/figure-pdf/affinity-1.pdf}

\begin{Shaded}
\begin{Highlighting}[]
\CommentTok{\# Check which words belong to each cluster}
\NormalTok{apck}
\end{Highlighting}
\end{Shaded}

\begin{verbatim}

APResult object

Number of samples     =  79 
Number of iterations  =  138 
Input preference      =  -0.1390897 
Sum of similarities   =  -2.322368 
Sum of preferences    =  -0.417269 
Net similarity        =  -2.739637 
Number of clusters    =  3 

Exemplars:
   can cant think
Clusters:
   Cluster 1, exemplar can:
      amazon bought can great love one reviews works
   Cluster 2, exemplar cant:
      cant doesnt good ipad like much price since work
   Cluster 3, exemplar think:
      able access also android anything app apps back battery better buy 
      computer day device dont easy email even fire first get got internet ive 
      just keep kindle little long looking many need never nice now product 
      purchased read really received recommend screen see something still store 
      support tablet take thing think time use used using want wanted way well 
      wifi will youre
\end{verbatim}

We notice that each cluster has an ``exemplar'' cluster, a prominent
word that best represents the cluster. It is a word that should define
the other words which are contained in the same cluster. It might be
considered the ``topic'' of the cluster and can be the thematic focus
of a cluster. In our case the exemplars are: \textit{can} , \textit{cant} and \textit{think}. Check
the words in each cluster and see if you can relate them to the
exemplar.

The examplars are chosen based on a message passing between data points
regarding their suitability as the most representative points of a
cluster. To demonstrate this, think of the reviews data set in our book. We want to
group reviews which are similar. Each review can ``send a
message'' to the other reviews revealing its topic. If another review has
the same topic, they will group. Some reviews will receive many messages
back, showing that they are popular within that topic. These
reviews become their groups' exemplars and have the strongest connections
to all the other reviews. The process continues until each review is
grouped. Eventually, the corpus will be organized into sections where
each group is characterized by a topic led by an exemplar that is the
distilled essence of the topic.

We might wonder which clustering algorithm performs optimally or yields
the most appropriate number of clusters for our specific data set.
Various methods and packages, like \texttt{cValid}, offer comparative
analyses of clustering algorithms by evaluating their stability and
validity across different data sets. However, the effectiveness of these
algorithms can be greatly dependent on the nature of the data set,
including its size, dimensions, distribution, and the presence of
outliers. Familiarity with these characteristics and hands-on experience
is crucial when selecting the most suitable clustering approach.
Practicing with diverse data sets and challenges is essential to
developing this expertise. For a comprehensive comparison of the
algorithms discussed here, including their strengths and limitations in
different scenarios, refer to Table A {[}reference{]} in the appendix.

\section{Document clustering}\label{documents-clustering}

In marketing, clustering similar elements enables more effective
segmentation. As an illustration, let us apply a method such
as HDBSCAN to cluster documents. Unlike in previous approaches, we will
use a document-term matrix (DTM) instead of a term-document matrix
(TDM), focusing on clustering documents with similar features. Although
we have processed the documents by removing elements like stop words and
numbers, this analysis aims to reveal insights based on the documents'
reduced set of features. For this purpose, we are using the same underlying
code but applying it to a DTM. This is achieved by transposing the matrix
so that rows and columns are switched, effectively turning the document
features into the primary axis for clustering. This allows us to examine
the thematic similarities between documents more directly.

\begin{Shaded}
\begin{Highlighting}[]
\FunctionTok{library}\NormalTok{(dbscan)}

\CommentTok{\# Matrix transposition}
\NormalTok{dtm }\OtherTok{\textless{}{-}} \FunctionTok{t}\NormalTok{(docs)}

\CommentTok{\# Calculate the hdbscan cluster}
\NormalTok{clustering.dbscan }\OtherTok{\textless{}{-}} \FunctionTok{hdbscan}\NormalTok{(dtm, }\AttributeTok{minPts =} \DecValTok{2}\NormalTok{)}

\CommentTok{\# Visualize the clusters created with hdbscan}
\FunctionTok{fviz\_cluster}\NormalTok{(}\FunctionTok{list}\NormalTok{(}\AttributeTok{data =} \FunctionTok{as.matrix}\NormalTok{(dtm), }\AttributeTok{cluster =}\NormalTok{ clustering.dbscan}\SpecialCharTok{$}\NormalTok{cluster), }
             \AttributeTok{ellipse.type =} \StringTok{"norm"}\NormalTok{,}
             \AttributeTok{repel =} \ConstantTok{TRUE}\NormalTok{)}
\end{Highlighting}
\end{Shaded}

\includegraphics{5.Chapter_files/figure-pdf/dbclusteringdocs-1.pdf}

\begin{Shaded}
\begin{Highlighting}[]
\CommentTok{\# Clustering details}
\NormalTok{clustering.dbscan}
\end{Highlighting}
\end{Shaded}

\begin{verbatim}
HDBSCAN clustering for 50 objects.
Parameters: minPts = 2
The clustering contains 2 cluster(s) and 45 noise points.

 0  1  2 
45  3  2 

Available fields: cluster, minPts, coredist, cluster_scores,
                  membership_prob, outlier_scores, hc
\end{verbatim}

In cluster 2, there are only two elements, documents 16 and 35. In R,
type \texttt{content{[}16{]}} and \texttt{content{[}35{]}} to compare
the documents.

The two reviews seem to be similar in terms of length and content. To
hone your skills, you can try various combinations of clustering
methods, distances, and parameters.

More generally, you can apply these techniques to any corpus in order to
extract their meaning in terms of topics, but also similarities and
differences between and within clusters.

\texttt{Exercise\ starts}

\textbf{Exercise 1: Hierarchical Clustering} Apply hierarchical
clustering using various distance metrics (e.g., Euclidean, Manhattan)
and linkage methods (e.g., single-linkage, average-linkage,
complete-linkage). Analyze how different combinations affect the
structure of the dendrogram.

\textbf{Exercise 2: K-means Clustering} Experiment with the k-means algorithm by
varying the number of clusters (K). Evaluate how the choice of K
influences the cohesion and separation of the clusters. Use methods like the
elbow or silhouette score to determine the optimal K value.

\textbf{Exercise 3: Gaussian Mixture Models (GMM) Clustering} Apply GMM
clustering to varying numbers of components (clusters). Experiment
with different types of covariance (e.g., full, tied, diagonal, spherical) and
initialization methods to observe their effects on the shapes and
sizes of clusters.

\textbf{Exercise 4: HDBSCAN Clustering} Explore the HDBSCAN algorithm by
modifying the minimum size of the clusters and the \(\epsilon\) value.
Experiment with adjusting the algorithm's parameters to see how they
affect the number of clusters detected and the handling of noise.

\textbf{Exercise 5: Affinity Propagation Clustering} Investigate the use
of different affinity propagation methods provided by the
\texttt{apcluster} package, such as \texttt{apcluster()} and
\texttt{apclusterL()}. Explore the impact of preference values and
similarity measures on the number and composition of the
exemplars and clusters.

\texttt{Exercise\ ends}

\textbf{Conclusion} \emph{Conclusion starts}

This chapter deals with the fundamentals of clustering analysis,
starting with basic step-by-step examples and advancing to more
sophisticated clustering methods. It covers a range of techniques,
including hierarchical clustering, K-means clustering, Gaussian mixture models
(GMM), DBSCAN, and affinity propagation, each selected to show its
unique approach to grouping data. The
chapter demonstrates how these methods can be applied to diverse
datasets using practical examples of coding, emphasizing the adaptability and utility of clustering in
uncovering hidden patterns within data. The discussion aims to provide
readers with a comprehensive understanding of how different clustering
algorithms perform. and their potential application to text analytics in
marketing.

\emph{Conclusion ends}
